\section{Discussion \& related work}

To the best of our knowledge, we didn't find any
research which explores the topic at hand, but the idea
is certainly in the air, as illustrated by the following
brilliant papers.

Original work by Lambek \cite{Lambek1980} sets the tone
of the future research, establishing correspondence
between type systems based on lambda calculus with
various kinds of categories, starting with
correspondence between simply typed lambda calculus and
cartesian closed categories.

Conal Elliott's ``Compiling to Categories''
\cite{Elliott:Compcat2017} shows how to interpret (some
subset of) Haskell in arbitrary \texttt{Category}
(again, definable in Haskell) via GHC compiler plugin.
That paper is a direct inspiration for this one. There
is an important difference, though: reinterpretation is
internal to our type systems, unlocking new patterns in
the usage of reinterpretation, including (but not
limited to) universe polymorphism, external
parametricity and more. However, Conal's approach is
more flexible since his translation to categories
requires desired category to only implement primitives
actually used for translation of a particular term at
hand.

In a recent paper ``A Certifying Proof Assistant for
Synthetic Mathematics in Lean'', Wojciech Nawrocki and
others \cite{Nawrocki:Lean2026} adapt Conal's approach
to Lean's system with dependent types, using
Categories-with-Families to represent internal models
and employing Lean's powerful metaprogramming facilities
instead of a compiler plugin. In doing so, they are able
to provide a framework for doing synthetic mathematics
in Lean, extracting proofs done in MLTT applied to more
specific categories. Thus our results in section
\ref{MLTT-R-synthetic} are close to the results shown in
their paper. To add to that, their work is already
implemented in practical proof assistant. Yet, our type
system involves no metaprogramming since
reinterpretation is internal to it, and it handles
conversion rule gracefully. In addition, the focus is a
bit different, where we use reinterpretation to express
other polymorphic features of a language.

Conor McBride's ``Crude but effective stratification''
\cite{McBride:Crude2002} (systemic treatment of this
idea is given by Favonia and others
\cite{Favonia:Displacement2023}) introduces an
eye-opening idea of using \textit{displacement} as a
lightweight form of universe polymorphism: to reuse term
typed at a lower universe level, lift all levels in a
term by adding a constant shift $s$. From the point of
view advocated by this paper, you can see it as an act
of reinterpreting an expression in an internal model of
type system where universe at level $i$ is interpreted
as a universe at level $i+s$ (section
\ref{MLTT-R-displacement}).

In ``Parametricity and Dependent Types'' by
Jean-Philippe Bernardy, Patrik Jansson and Ross Paterson
\cite{Bernardy:Parametricity2010}, authors introduce a
direct translation procedure to yield parametricity
theorems and proofs from polymorphic definitions in a
dependently typed system. The definition of <<external
parametricity>> we follow, though, is given in
``Internal Parametricity, without an Interval'' by
Thorsten Altenkirch, Yorgo Chamoun, Ambrus Kaposi and
Michael Shulman \cite{Altenkirch:Parametricity2024}.
Once again, we can reformulate the translation they give
as reinterpretation in an internal model. Altenkirch and
others noted an issue with this approach: if we apply
the external translation to the judgment with nonempty
context, context becomes parametric, too; that's the
reason they add \texttt{rel} operation and yield type
system with \textit{computational properties of higher
observational type theory}. We sidestep the issue with
context since reinterpretation is applied only to terms
whose context is also <<reinterpretable>> and is not
affected by the translation (section
\ref{MLTT-R-parametricity}).
